\chapter{Artes visuales y construcción de identidad}
Wind muestra cómo los artistas renacentistas reintrodujeron símbolos paganos para comunicar virtudes cívicas y tensiones teológicas, reconfigurando el imaginario cristiano desde la mitología grecolatina.\cite{Wind1958} Ames-Lewis destaca que esta apropiación iconográfica exigía una cultura visual sofisticada, capaz de dialogar con textos antiguos y con los mecenas que encargaban las obras.\cite{AmesLewis2000}

La autorrepresentación ocupa un lugar central en el periodo. Woods-Marsden analiza el autorretrato como afirmación del estatuto social del artista, quien pasa de ser un artesano a un intelectual que negocia su imagen pública.\cite{WoodsMarsden1998} Este cambio se refleja en los formatos de los talleres, en la circulación de bocetos y en la presencia del pintor dentro de la escena representada.

\begin{figure}[H]
    \centering
    % Reemplazar por la imagen recortada correspondiente
    \includegraphics[width=0.7\textwidth]{FIGURES/autorretrato_placeholder}
    \caption{Autorretrato con herramientas de taller ("Joanna Woods-Marsden - Renaissance Self-Portraiture: The Visual Construction of Identity and the Social Status of the Artist (1998, Yale University Press) - libgen.li.pdf", p. XX).}
    \label{fig:autorretrato}
\end{figure}

La práctica del dibujo preparatorio permitió codificar gestos, proporciones y emociones, generando un lenguaje visual compartido. Las academias y los tratados teóricos ampliaron ese vocabulario, consolidando la figura del artista como intérprete erudito del pasado clásico.
